% SGA 3，报告 VII_B，英文重建本——第 1 轮（中文译本）。
% 范围：依据原文件第 14--15 页，完整包含 0.4 至 0.4.2；止于原文件第 16 页的 0.5 之前。
% 合并阅读版页码： 527--528. 再版印刷页码： 516--517.
% 本范围内的原版页边定位： 488--490.
% 依据： Polo--Gille Exp7B-13oct24.pdf,
% SHA-256 D303690E0195654BBE7C3DAA5D89EA23E24CD709A3575217234AC19AE8836694.
% 曾参阅 Jacob Reinhold 的 CC BY 4.0 Markdown，作为对比版本来源。

\section*{0.4. 预有限代数}
\phantomsection
\label{sga3:VIIB:0.4}
\addcontentsline{toc}{subsection}{\SGARefLink{sga3:VIIB:0.4}{0.4}. 预有限代数}

设 \(k\) 是伪紧环。一个\emph{拓扑 \(k\)-代数}是交换拓扑环 \(A\)，
并配备一个拓扑环态射 \(k\to A\)。若这个代数的底拓扑
\(k\)-模是伪紧的，则称 \(A\) 为\emph{预有限 \(k\)-代数}。

设 \(\mathfrak l\) 是这样的 \(A\) 的开 \(k\)-子模。复合映射
\[
  \varphi:
  A\times A
  \xrightarrow{\mathrm{mult}}A
  \longrightarrow A/\mathfrak l
\]
是连续的。因此，\SGARefLink{sga3:VIIB:0.3.1}{引理 0.3.1} 给出 \(A\) 的开 \(k\)-子模
\(\mathfrak n\)，使 \(\varphi(A\times\mathfrak n)=0\)。所以
\(\mathfrak l\) 包含开理想 \(A\mathfrak n\)，且 \(A\) 是伪紧环。

同样，设 \(M\) 是底 \(k\)-模伪紧的拓扑 \(A\)-模。若 \(M'\) 是
\(M\) 的开 \(k\)-子模，对
\[
  A\times M
  \xrightarrow{\mathrm{mult}}M
  \longrightarrow M/M'
\]
应用\SGARefLink{sga3:VIIB:0.3.1}{引理 0.3.1}，表明 \(M'\) 包含开 \(A\)-子模。
因此，\(M\) 也是伪紧 \(A\)-模。\footnote{编者添加了下列引理；参见\
\SGARefLink{sga3:VIIB:0.5}{命题 0.5} 的编者注~(36)。}
反过来：

\medskip
\noindent\textbf{引理。}
\phantomsection
\label{sga3:VIIB:0.4-restriction-lemma}
\textit{设 \(A\) 是预有限 \(k\)-代数，\(M\) 是伪紧 \(A\)-模。
通过限制标量得到的 \(k\)-模 \(M|_k\) 是伪紧的。}

事实上，每个有限长度的伪紧 \(A\)-模都同构于某个商
\[
  A^n/L,
\]
其中 \(L\) 是 \(A^n\) 的开子模，因此该商是伪紧 \(k\)-模。
由于 \(M|_k\) 是这种模的投射极限，它是伪紧 \(k\)-模。

\medskip
\noindent\textbf{0.4.1.}
\phantomsection
\label{sga3:VIIB:0.4.1}
若 \(A\) 和 \(B\) 是预有限 \(k\)-代数，则从 \(A\) 到 \(B\) 的
\emph{态射}按定义是连续的 \(k\)-代数同态。我们将预有限 \(k\)-代数范畴记为
\[
  \operatorname{Alp}/k.
\]
该范畴显然有投射极限：投射极限的底代数是各底代数的投射极限，赋予投射极限拓扑。

它也有有限直极限。\footnote{在 \SGARefLink{sga3:VIIB:0.4.2}{0.4.2} 中将看到，
它有任意直极限。}例如，若
\[
  f:A\longrightarrow B,
  \qquad
  g:A\longrightarrow C
\]
是预有限 \(k\)-代数的态射，则 \(B\) 与 \(C\) 在 \(A\) 上的融合和的底拓扑 \(A\)-模为
\[
  B\widehat{\otimes}_A C
\]
这里，\SGARefLink{sga3:VIIB:0.4}{0.4} 通过 \(f\) 和 \(g\) 使 \(B\) 和
\(C\) 成为伪紧 \(A\)-模，而乘法由下式刻画
\[
  (b\widehat{\otimes}c)
  (b'\widehat{\otimes}c')
  =
  bb'\widehat{\otimes}cc'
\]
其中 \(b,b'\in B\)，\(c,c'\in C\)。

% 依据： Exp7B 原文件页码 15; 合并阅读版页码 528; 印刷页码 517.
\medskip
\noindent\textbf{定义 0.4.2。}
\phantomsection
\label{sga3:VIIB:0.4.2}
\textit{若预有限 \(k\)-代数 \(C\) 的底 \(k\)-模长度有限（因而是离散的），
则称其具有\emph{有限长度}。以 \(\operatorname{Alf}/k\) 表示
\(\operatorname{Alp}/k\) 中由有限长度的 \(k\)-代数组成的满子范畴。}\footnote{注意：
只有当 \(k\) 是 Artin 环时，\(k\) 才是 \(\operatorname{Alf}/k\) 的对象；
参见下文 \SGARefLink{sga3:VIIB:1.2.2}{1.2.2}。}

对每个预有限 \(k\)-代数 \(A\)，令 \(h_A\) 为函子
\[
  h_A:\operatorname{Alf}/k\longrightarrow(\mathbf{Sets}),
  \qquad
  C\longmapsto
  \operatorname{Hom}_{\operatorname{Alp}/k}(A,C).
\]
函子 \(h_A\) 左正合，\footnote{即它与有限投射极限交换。}
且典范投影 \(A\to A/\mathfrak l\)（其中 \(\mathfrak l\) 遍历 \(A\) 的开理想）
对每个 \(C\in\operatorname{Alf}/k\) 诱导关于 \(C\) 函子性的同构
\[
  \varinjlim_{\mathfrak l}
  \operatorname{Hom}_{\operatorname{Alf}/k}(A/\mathfrak l,C)
  \xrightarrow{\sim}
  \operatorname{Hom}_{\operatorname{Alp}/k}(A,C),
\]
换言之，\(h_A\) 是可表函子 \(h_{A/\mathfrak l}\) 的直极限：
\begin{equation}
  h_A\simeq\varinjlim_{\mathfrak l}h_{A/\mathfrak l}.
  \tag{\(\ast\)}
  \label{sga3:VIIB:0.4.2-star}
\end{equation}

若 \(B\) 是另一个预有限 \(k\)-代数，则由双函子 \(\operatorname{Hom}\)
的一般性质，以及对有限长度的 \(C\) 成立的等式
\[
  \operatorname{Hom}(h_C,h_B)=h_B(C)
\]
得到同构
\begin{align*}
  \operatorname{Hom}_{\operatorname{Alp}/k}(B,A)
  &\simeq
  \varprojlim_{\mathfrak l}
  \operatorname{Hom}_{\operatorname{Alp}/k}(B,A/\mathfrak l)
  \\
  &\simeq
  \varprojlim_{\mathfrak l}
  \operatorname{Hom}(h_{A/\mathfrak l},h_B)
  \\
  &\simeq
  \operatorname{Hom}
    \left(\varinjlim_{\mathfrak l}h_{A/\mathfrak l},h_B\right).
\end{align*}
结合 \eqref{sga3:VIIB:0.4.2-star}，这证明反变函子 \(A\mapsto h_A\) 是全忠实的。
事实上：

\medskip
\noindent\textbf{命题。}
\phantomsection
\label{sga3:VIIB:0.4.2-proposition}
\textit{函子 \(A\mapsto h_A\) 是从 \(\operatorname{Alp}/k\) 到从
\(\operatorname{Alf}/k\) 至 \((\mathbf{Sets})\) 的左正合函子范畴的反等价。}

根据上文，只需证明每个左正合函子
\[
  F:\operatorname{Alf}/k\longrightarrow(\mathbf{Sets})
\]
都同构于某个函子 \(h_A\)。代数 \(A\) 可以如下构造（参见\ TDTE~II，\S~3）。

由于 \(F\) 左正合，对每个有限长度的 \(k\)-代数 \(C\) 及每个
\(\xi\in F(C)\)，存在 \(C\) 的最小子代数 \(C'\)，使 \(\xi\) 属于
\(F(C')\) 在 \(F(C)\) 中的像。若 \(C'=C\)，则称对 \((C,\xi)\) 为\emph{极小}的。

将从 \((C,\xi)\) 到 \((D,\eta)\) 的态射定义为满足
\[
  (F\phi)(\xi)=\eta.
\]
的同态 \(\phi:C\to D\)，极小对就构成一个范畴。
这样的 \(\phi\) 必为满射，且极小对范畴是左滤过的。此外，可以只考虑那些 \(C\)
属于一个包含每种同构类型的一个有限长度 \(k\)-代数的集合的对 \((C,\xi)\)。\footnote{
事实上，每个长度为 \(n\) 的离散 \(k\)-模都同构于 \(k^n/L\)，其中
\(L\) 是 \(k^n\) 中余长度为 \(n\) 的开子模。然后可以考虑每个这样的模上
预有限 \(k\)-代数结构的同构类所成的集合。}
因此，从极小对范畴到预有限 \(k\)-代数范畴的函子 \((C,\xi)\mapsto C\)
有投射极限；取 \(A\) 为此极限。

\medskip
\noindent\textbf{推论。}
\phantomsection
\label{sga3:VIIB:0.4.2-corollary}
\textit{范畴 \(\operatorname{Alp}/k\) 有任意直极限。}

事实上，从 \(\operatorname{Alf}/k\) 到 \((\mathbf{Sets})\) 的左正合函子范畴
有逐点计算的投射极限：若 \((F_i)\) 是这类函子的投射系，则对每个
\(C\in\operatorname{Alf}/k\)，
\[
  (\varprojlim_iF_i)(C)=\varprojlim_iF_i(C).
\]
\footnote{若 \((A_i)_{i\in I}\) 是 \(\operatorname{Alp}/k\) 中的直系，则其在
\(\operatorname{Alp}/k\) 中的直极限是底 \(k\)-代数的直极限 \(k\)-代数 \(B\)
对于由如下理想 \(I\) 定义的拓扑的分离完备化：对每个 \(i\)，\(I\cap A_i\)
在 \(A_i\) 中开，且 \(\operatorname{length}_k(B/I)<\infty\)。例如应注意，
若 \(K/k\) 是无限次数的域代数扩张，\((k_i)\) 是有限次子扩张的滤过直系，
则系 \((k_i)\) 在 \(\operatorname{Alp}/k\) 中的直极限是零环。}

% 下一断点： Exp7B 原文件页码 16, 合并阅读版页码 529,
% 再版印刷页码 518, 0.5 处。
